%% --------------------------------------------------------
\chapter{Поляризація CMB}
\label{sec:pol}
%% --------------------------------------------------------

Температурний спектр $C_\ell^{TT}$, розглянутий у попередньому
розділі, несе інформацію переважно про скалярні збурення.
Але CMB містить і другий канал інформації~--- \emph{поляризацію}.
Вона виникає при томсонівському розсіянні фотонів і дозволяє
розрізнити два типи первинних збурень: скалярні (флуктуації
густини) і тензорні (гравітаційні хвилі). Детекція тензорного
сигналу у поляризації CMB стала б прямим підтвердженням
квантової природи гравітації~--- саме тому пошук B-мод є
однією з центральних задач сучасної космології.

%% --------------------------------------------------------
\section{Механізм поляризації: розсіяння Томсона}
%% --------------------------------------------------------

Фотони поляризуються при останньому розсіянні на електронах.
Томсонівське розсіяння поляризує фотон \emph{лише} якщо навколо електрона
існує \textbf{квадрупольна} температурна анізотропія: якщо температура
однакова з усіх боків~--- поляризації немає.

Безпосередньо до рекомбінації фотонів і електронів ще достатньо для
ізотропізації, тому квадруполь малий і поляризація слабка.
Але під час самої рекомбінації ($\Delta z \sim 80$) оптична глибина
падає швидко~--- квадруполь встигає вирости і поляризація «відбивається»
в поверхні останнього розсіяння.

Напрямок поляризації залежить від \textbf{поперечного} перепаду температури
навколо точки розсіяння.

%% --------------------------------------------------------
\section{E- і B-моди}
%% --------------------------------------------------------

Поляризаційне поле на сфері, як і будь-яке тензорне поле,
однозначно розкладається на дві компоненти~--- аналогічно
до розкладу вектора на потенціальну і вихрову частини.
Ці дві компоненти мають різне фізичне джерело і різний
статус спостережень:

\begin{center}
\begin{tblr}{XXX}
\toprule
Тип & Джерело & Статус \\
\midrule
E-моди (градієнтний, симетричний візерунок) & Скалярні флуктуації & Виміряні \\
B-моди (ротаційний, закручений візерунок)  & Тензорні флуктуації (GW) & Поки не знайдені \\
\bottomrule
\end{tblr}
\end{center}

Чому два типи збурень дають різні моди поляризації~---
питання симетрії, а не лише статистики:

\begin{keybox}
\textbf{Чому скалярні збурення не дають B-мод?}
Скалярні збурення мають потенціальний (паритетно-парний) характер:
$\Phi(\mathbf{x})$ визначає поле швидкостей $\mathbf{v} = \nabla\Phi$.
E-моди мають ту ж симетрію паритету~--- вони парні.
B-моди псевдоскалярні (непарні), і лінійне скалярне збурення їх не генерує.
Тензорне збурення $h_{ij}$ має обидва стани, в тому числі непарний~---
звідси B-моди від гравітаційних хвиль.
\textit{Практичний забруднювач:} гравітаційне лінзування на великих масштабах перетворює
E-моди в B-моди~--- це треба враховувати при пошуку первинних B-мод.
\end{keybox}

Знайти B-моди~--- означає отримати прямий доказ первинних гравітаційних хвиль
і квантової природи гравітації в режимі малих збурень.
Поточна верхня межа: $r < 0.036$ (95\% CL, BICEP/Keck~2021 \cite{bicep_keck_2021}).

Таким чином, CMB є повноцінним «двоканальним» джерелом інформації:
температурний спектр обмежує скалярні збурення і параметри $\Lambda$CDM,
а поляризаційні B-моди~--- єдиний прямий зонд тензорних збурень і
енергетичного масштабу інфляції. Обидва канали разом утворюють
найпотужніший інструмент для перевірки інфляційних моделей~---
включно з тією, що розглядається у наступних розділах.
